%% supplementary.tex --- Supplementary material for the IST submission
%% Compile standalone with pdfLaTeX. Not counted toward the journal's word limit.
\documentclass[11pt]{article}
\usepackage[T1]{fontenc}
\usepackage{lmodern}
\usepackage[margin=1in]{geometry}
\usepackage{booktabs}
\usepackage{url}
\usepackage[hidelinks]{hyperref}

\title{Supplementary material for\\
\emph{Typed Issue Specifications from App Reviews: What Scales, and a Downstream Benefit That
Does Not Replicate}}
\author{Fabiha Jalal \and Sadia Tasnim Dhruba \and Tazrian Zaman Tushti \and Ajwad Abrar \and
Md Kamrul Hasan \and Hasan Mahmud}
\date{}

\begin{document}
\maketitle

This file carries the implementation settings and the annotator instruments. It is referenced
from the manuscript and is not required to follow the argument there. We separate it because an
audit of the manuscript against its own text found that every instrument was named but not
specified, so a reader could reconstruct the architecture and not a single number. What follows
is the repair, at a level of detail the manuscript has no room for.

\section{Implementation Settings}\label{appendix:settings}

  Section~\ref{system-design} states what each stage does and why. This appendix gives the
  settings. An audit of this paper against its own text found that every instrument was named but
  not specified, so a reader could reconstruct the architecture and not a single number. What
  follows is the repair. Values are as they appear in the released code; where a setting is a
  default we say so.

  \textbf{Stage 1, classification.} Backbone \texttt{roberta-base} (cased, 125M), multi-label
  sigmoid head over seven classes, per-label cut 0.5. Fine-tuning: 3 epochs, learning rate
  $2\times10^{-5}$, batch size 16, max sequence length 256, AdamW at library defaults, seed 42.
  Training corpus is the relabelled corpus under a stratified per-class cap of 15{,}000 rows,
  shuffled at seed 42; class order follows Python dictionary insertion order, which
  Section~\ref{exp1-results} shows is load-bearing for the held-out split. HITL-1 routes a review
  to a human when top-class probability $< 0.70$ or the top-two margin $< 0.15$.

  \textbf{Stage 1, the V2 annotator.} The noisy corpus labels come from
  \texttt{gemini-2.0-flash} at temperature 0.1, 150 max output tokens, prompted for a single
  primary class plus a self-reported confidence, with four written tie-break rules (device-specific
  faults route to \texttt{compatibility}, ``slow'' and ``lag'' to \texttt{performance}, and so on).
  The prompt is released in full. On any API error or parse failure the caller returns
  \texttt{other} with confidence 0.0, so a reader cannot distinguish a genuine \texttt{other} from
  a failed call; this is a defect we report rather than a design choice. V1, V3 and V4 are
  intermediate checkpoints kept for the record and carry no reported result.

  \textbf{Stage 1, the verified-anchor procedure.} Because this is the contribution we defend, we
  state it as an algorithm. (i) The anchor is 5{,}230 reviews the lead author labelled by hand,
  drawn from the strata where the V2 labeller was least productive; 5{,}038 of them are
  \texttt{praise}, so the anchor is not a balanced sample and we do not claim it is one. (ii)
  Out-of-sample predicted probabilities come from \texttt{sklearn.cross\_val\_predict} with
  5 folds over a TF-IDF plus linear pipeline. (iii) \texttt{cleanlab.filter.find\_label\_issues}
  is called with those probabilities and the noisy labels, ranked by self-confidence, and
  \texttt{get\_label\_quality\_scores} supplies a per-row score. (iv) A flagged row is corrected to
  the anchor's label only when the anchor's confidence is $\geq 0.70$ and the V2 label's
  probability is $\leq 0.20$; otherwise the original label is kept. (v) Rows the anchor does not
  cover, which is 97.6\% of the corpus, are corrected to cleanlab's argmax under the same
  probability ceiling. The $3\times3$ grid over $(0.65, 0.70, 0.75) \times (0.15, 0.20, 0.25)$ in
  Section~\ref{exp1-results} sweeps (iv). ``Anchor confidence'' is the annotator's recorded
  confidence in their own label, not a model quantity.

  \textbf{Stage 1, compatibility augmentation.} 200 template-generated synthetic reviews (device,
  OS, and screen templates) plus 100 keyword-mined RRGen reviews. Section~\ref{exp1-results}
  reports two defects in this augmentation that a reader needs before using the class.

  \textbf{Stage 2.} Embeddings from \texttt{all-MiniLM-L6-v2}. UMAP with
  \texttt{n\_components}~=~50, \texttt{min\_dist}~=~0.0, \texttt{metric}~=~cosine, and
  \texttt{n\_neighbors} tuned per class (30 for \texttt{bug\_report}, 25 for
  \texttt{feature\_request}, 20 for \texttt{performance} and \texttt{usability}, 10 for
  \texttt{compatibility}). HDBSCAN on the reduced space with Euclidean metric and per-class
  \texttt{min\_cluster\_size} / \texttt{min\_samples} of 200/20, 100/10, 60/8, 60/8, and 10/3
  respectively. Aspects are spaCy noun chunks filtered against a curated keyword list, released
  with the code. PageRank runs undirected on the bipartite review-aspect graph at NetworkX
  defaults (damping 0.85). The KG-hierarchical side runs over an 8{,}404-review sample; the flat
  baseline runs over 91{,}368 clustered reviews from the full corpus, so the two describe
  different review sets (Section~\ref{results}).

  \textbf{Stage 3.} Headline generator \texttt{claude-opus-4-7}; comparators
  \texttt{llama-3.3-70b-versatile} via Groq at temperature 0.2 and 1500 max tokens, and
  \texttt{Qwen/Qwen2.5-3B-Instruct} and \texttt{Qwen/Qwen2.5-1.5B-Instruct} decoded greedily at
  600 max tokens. None of these identifiers is a dated snapshot, which is a reproducibility defect
  we should have avoided. The generation prompt is released in full and is byte-identical across
  the local comparators. Routing uses the V5 predicted class of the cluster; \texttt{praise} and
  \texttt{other} have no template and fall back to the bug-report form, which means the pipeline
  emits no meaningful specification for the corpus's largest class. The pipeline re-prompts at
  most twice with dimension-level feedback when the mean rubric score is below 3.0.

  \textbf{Stage 4.} ChromaDB over 15{,}100 documents from five sources (10{,}000 RRGen pairs,
  5{,}000 similar responses, 60 changelogs, 40 FAQ entries, and the IssueSpec injected per query
  rather than pre-indexed). The re-run generator is \texttt{claude-opus-4-7} in both arms. Neither
  arm's file stores its prompt or its retrieved context, so the manipulation is not recoverable
  from the artifact; Section~\ref{exp2-results} states what follows from that.

  \textbf{Stage 5.} Five policies on distilGPT2 (82M) from 400 SFT samples, 296 KTO pairs, and 100
  DPO pairs; the Lagrangian loop runs 30 steps at batch size 4 with KL coefficient 0.05, and the
  compliance sweep covers $\tau \in \{0.0, 0.5, 0.7, 0.9, 1.0\}$. The binding-constraint re-run
  uses \texttt{Qwen/Qwen2.5-1.5B-Instruct} with LoRA ($r$~=~8, $\alpha$~=~16, on
  \texttt{q,k,v,o\_proj}) for 90 steps.

  \textbf{Instruments.} Strict template-fill requires, per field: title $\geq 4$ words;
  description $\geq 30$ words; $\geq 3$ reproduction steps, each $\geq 5$ words and containing one
  of a 33-verb action list released with the code; expected and actual behaviour $\geq 8$ words
  each; $\geq 3$ acceptance criteria of $\geq 8$ words; a severity in P0--P3; and an affected
  component of $\geq 2$ words not on a 10-item generic stoplist. All comparisons are
  ``at least''. \textsc{SpecCov} counts substantive-token overlap with the source cluster, tokens
  of length $\geq 5$ after stop-word filtering, with quoted-string and coverage bonuses, banded to
  1--5; Section~\ref{speccov-validation} withdraws it as a faithfulness measure. The
  five-dimension rubric and the Y/P/N purity rubric are scored by
  \texttt{Qwen/Qwen2.5-3B-Instruct} decoded greedily; both judge prompts are released, the judge
  does not see the source reviews when scoring the accuracy dimension, and a spec whose output
  does not parse into all five dimensions is dropped without a counter. The compliance scorer is a
  regular expression over roughly thirty tokens across four dimensions; Section~\ref{exp3-results}
  reports what it does and does not detect.

\section{Annotator Guidelines}\label{appendix:annotators}

  The lead author (PhD in software engineering with prior mobile-app review-mining experience)
  served as primary annotator across all five annotation tasks. The 490-review classification
  gold was additionally annotated by two further independent human raters on the identical item
  set, aligned by review id; all three first completed a shared 20-review calibration round,
  disagreements were adjudicated before the main task opened, and the released gold label is the
  majority vote over the three raters. All human annotation in this work is carried out by
  those same three people. For Stage~4, all three raters scored all 200 rows of the reported
  re-run independently and blind to condition; the discarded 400-row round had the lead author as
  primary rater with one additional rater on a 30-pair subset, and is described here only because
  the paper reports its withdrawal. Separately, two Qwen2.5-3B-Instruct prompts (concise
  role-based and chain-of-thought) served as the second and third raters in the 99-review
  cluster-purity Krippendorff $\alpha$ calibration, which is an LLM-assisted panel and not one of
  the human rounds. For the 5{,}230-record verified anchor, which covers 5{,}229 distinct corpus rows, annotators classified each review into
  one of the seven classes in our extension of the Maalej-Nabil taxonomy (\texttt{bug\_report}, \texttt{feature\_request},
  \texttt{performance}, \texttt{usability}, \texttt{compatibility}, \texttt{praise},
  \texttt{other}) with multi-label assignment permitted when a review carried multiple concerns;
   reviews below the ``fairly confident'' threshold were skipped and excluded from the anchor.
  For the 490-review classification gold standard, a stratified sample of 70 reviews per class
  was drawn from the working corpus, and the lead author re-labeled each review independently
  and blinded to the V2 LLM prediction; the 490 gold labels serve as the held-out evaluation set
   for Cohen's $\kappa$ progression (Table~\ref{tab:kappa}) and per-class F1
  (Table~\ref{tab:perclass-f1}). For the 5-dimension IssueSpec rubric (all 100 specs; an earlier stratified $n=28$ subsample, 6 per issue
   type plus 2 extras), each spec was scored on a 1--5 Likert scale across five dimensions:
  \emph{completeness} (are all required template fields filled?), \emph{accuracy} (do field
  values correctly reflect the source cluster?), \emph{actionability} (would a developer know
  what to do?), \emph{specificity} (does the spec name concrete entities rather than generic
  terms?), and \emph{clarity} (is the spec easy to read and unambiguous?), with each dimension
  scored independently. For the reported Stage~4 evaluation (100 reviews $\times$ 2 conditions,
  the same model with and without the IssueSpec, 200 rows), all three raters rated
  \emph{quality} (1--5), \emph{specificity} (1--5), and \emph{helpful} (Y/N) while blinded to
  the generation method; the superseded 400-row round used four template-composed conditions and
  is retained in the release for comparison. In both rounds condition labels were replaced with random A/B/C/D codes and unsealed
  only at aggregation, with the assignment persisted as spreadsheets with JSON master keys for
  audit. A second human rater independently scored a 30-pair (\texttt{no\_spec}, \texttt{full})
  subset for inter-rater calibration. For the Y/P/N cluster-purity rubric (initial 50-cluster
  audit; 99-cluster three-rater subsample), annotators read 5 representative reviews per cluster
   and assigned \emph{Y} (all 5 share the same sub-theme within the cluster's aspect), \emph{P}
  (3 or 4 of 5 share the sub-theme), or \emph{N} (fewer than 3 share); weighted purity is
  $\mathrm{purity}_w = (|Y| + 0.5|P|)/(|Y|+|P|+|N|)$, with $Y$ contributing 1.0, $P$
  contributing 0.5, and $N$ contributing 0. All comparative evaluations were blinded throughout
  the rating window and unsealed only at aggregation; reviews or specs below the annotator
  confidence threshold were marked as ``skip'' and excluded from downstream analysis.

\end{document}
